problem:  
Let \(\mathrm{Imm}^{2}(S^{1},\mathbb{R}^{2})\) be the Banach manifold of \(H^{2}\)-immersed closed planar curves and \(\mathcal{S} = \mathrm{Imm}^{2}/\mathrm{Diff}^{2}(S^{1})\) the corresponding shape space.  
For fixed \(A,B>0\), consider the second‑order Sobolev Riemannian metric  
\[
G^{(2)}_c(h,k)=\int_{S^{1}}\big(\langle h,k\rangle + A\langle D_s h, D_s k\rangle + B\langle D_s^{2}h, D_s^{2}k\rangle\big)\,ds,
\]
where \(D_s = \partial_{\theta}/|c'(\theta)|\).  
The associated **geodesic spray** \(F:T\mathrm{Imm}^{2}\to TT\mathrm{Imm}^{2}\) defines a smooth vector field whose integral curves are the geodesics of \(G^{(2)}\).

---

**Problem (analyticity of the Sobolev \(H^{2}\) geodesic spray and flow):**

1. (**Analyticity of the spray**)  
   Is the geodesic spray \(F\) **real‑analytic** as a map between Banach manifolds modeled on \(H^{2}(S^{1},\mathbb{R}^{2})\)?  
   That is, for each point \((c,v)\in T\mathrm{Imm}^{2}\), does there exist an open neighborhood \(U\) such that \(F|_U\) has a convergent power‑series representation in \((c,v)\), or equivalently admits locally bounded derivatives of all orders satisfying the usual factorial‑growth analyticity bounds in the Banach‑analytic sense?

2. (**Analytic dependence of the geodesic flow**)  
   Assuming \(F\) is analytic, does the local geodesic flow map
   \[
   \Phi:(t,c_0,v_0)\mapsto (c(t;c_0,v_0),\partial_t c(t;c_0,v_0))
   \]
   depend **holomorphically** on the (complexified) time variable \(t\) near \(t=0\) for fixed \((c_0,v_0)\)?  
   In other words, can the flow be extended to a holomorphic map on a strip \(\{t\in\mathbb{C}:|\mathrm{Im}(t)|<\varepsilon\}\subset\mathbb{C}\) taking values in \(T\mathrm{Imm}^{2}\)?

---

Why is it a "good" problem:  
1. **Profound insight:**  
   This problem clarifies whether the intrinsic dynamics of the \(H^{2}\) Sobolev metric on shape space are not only smooth but genuinely analytic. Establishing analyticity would reveal a new level of regularity and structural rigidity in infinite‑dimensional Riemannian geometry, comparable to analytic continuation phenomena in finite‑dimensional symmetric spaces.  

2. **Bridge between fields:**  
   It joins *geometric analysis* and *Banach‑space analytic PDE theory*, directly importing ideas from analytic semigroups and holomorphic ODE flows into the context of nonlinear differential geometry on spaces of embeddings. Such a connection could enable new analytical tools for studying shape evolution, curvature control, and stability.  

3. **Extreme simplicity:**  
   The essential question is elegantly simple:  
   *“Is the \(H^{2}\) elastic geodesic spray analytic, and can its flow be extended holomorphically in time?”*  
   Despite the clear formulation, resolving it would require developing a fine‑grained analytic calculus for nonlinear arc‑length‑dependent operators on Sobolev manifolds—representing the archetypal “narrow entrance, profound depth’’ challenge of a first‑rate research problem.